% Double-anonymous evidence-gated protocol/preprint for an Elsevier target.
% Do not input the separate, non-anonymous title_page.tex in this review file.
% Keep \resultslockedtrue until the sealed confirmatory run is complete.
% VALIDATION_MARKERS: RESULTS_LOCKED COMMIT CLARIFY ABSTAIN
\documentclass[preprint,12pt,authoryear]{elsarticle}

\usepackage[T1]{fontenc}
\usepackage[utf8]{inputenc}
\usepackage{amsmath,amssymb}
\usepackage{booktabs}
\usepackage{tabularx}
\usepackage{array}
\usepackage{graphicx}
\usepackage{xcolor}
\usepackage{url}
\usepackage{hyperref}
\usepackage{enumitem}

\biboptions{round,authoryear}
\journal{Engineering Applications of Artificial Intelligence}

\definecolor{lockred}{RGB}{145,30,30}
\definecolor{softgray}{RGB}{245,245,245}

\newif\ifresultslocked
\resultslockedtrue

\newcommand{\cdcv}{\textsc{CDCV-Gate}}
\newcommand{\commit}{\ifmmode\operatorname{COMMIT}\else\textsc{Commit}\fi}
\newcommand{\clarify}{\ifmmode\operatorname{CLARIFY}\else\textsc{Clarify}\fi}
\newcommand{\abstain}{\ifmmode\operatorname{ABSTAIN}\else\textsc{Abstain/Escalate}\fi}
\newcommand{\iedi}{\textsc{IEDI}}
\newcommand{\resultcell}{\ifresultslocked\texttt{RESULTS\_LOCKED}\else
  \textbf{INSERT VERIFIED RESULT}\fi}
\newcommand{\resultlockbox}{%
  \begin{center}
  \fcolorbox{lockred}{softgray}{%
    \begin{minipage}{0.92\linewidth}\small
    \textcolor{lockred}{\textbf{RESULTS LOCKED.}}
    This manuscript specifies a confirmatory study; it does not yet report a
    completed empirical evaluation. Every result cell remains locked until a
    community-validated, family-disjoint test set has been sealed, all
    calibration decisions have been frozen, and the prespecified analysis has
    been run. Synthetic design simulations are not empirical evidence.
    \end{minipage}}
  \end{center}}

\begin{document}

\begin{frontmatter}

\title{When to Interpret and When to Ask: Counterfactual Dialect-Context
Verification for Selective Intended-Meaning Resolution\tnoteref{tn-verification}}

\tnotetext[tn-verification]{Protocol/preprint. Here ``verification'' denotes
behavioral probing of prediction stability and candidate-relative
responsiveness, together with schema, provenance-category, and conflict
checks. It does not authenticate dialect identity, establish a speaker's
private intention, or certify contextual truth.}

\author{Anonymous for review}

\begin{abstract}
An utterance can be transcribed correctly yet interpreted incorrectly when its
intended interactional meaning depends on local conventions, discourse goals,
or a preceding speech act. Supplying a language model with more context does
not by itself establish that the context is sufficient, coherent, or being
used appropriately. We present the design of the Counterfactual
Dialect-Context Verification and Clarification Gate (\cdcv), a
schema-constrained inference-time mechanism for selective intended-meaning
resolution. The mechanism scores fixed candidate meanings using an ephemeral,
episode-specific context card; probes the model-selected candidate using
instances to be community validated as hypothesis preserving or hypothesis
contrasting;
and aggregates confidence, candidate-conditioned invariance, targeted
response, context-conflict, and validity signals. A behavioral-consistency
controller calibrated on
expression-family-disjoint data then commits to a sense, requests one
validated context-slot clarification, or abstains and optionally escalates.
We prespecify a benchmark design with 150 intended sealed-test families
across English use in the United States, Nigeria, and Korea, with sufficient,
alternate, underspecified, and
contradictory action labels. The primary estimand is wrong-commit risk at
matched coverage against an equal-call/token structured-context control;
secondary outcomes quantify invariance, appropriate counterfactual response,
clarification repair, per-community behavior, and computation. Current IEDID
rows and inherited personas are restricted to discovery, development, schema
construction, and lineage diagnostics. They are excluded from sealed-test
gold. A separately labelled, consented text-prototype pilot is prespecified to
assess bounded application feasibility after all models and thresholds are
frozen. This protocol/preprint makes no performance or deployment claim; it
defines the mechanism, evaluation contract, and evidence gates required for
one. The gate does not verify semantic truth, ground-truth intended meaning,
or context authenticity.
\end{abstract}

\begin{keyword}
selective prediction \sep pragmatic ambiguity \sep World Englishes \sep
counterfactual testing \sep clarification \sep context verification
\end{keyword}

\end{frontmatter}

\resultlockbox

\section{Introduction}

Language systems can capture every word and still miss what a speaker is
doing with those words. Pragmatic interpretation depends on linguistic
co-text, common ground, discourse goals, activity type, and locally meaningful
contextualization cues \citep{grice1975logic,gumperz1982discourse,
clark1996using,frank2012predicting}. The resulting ambiguity is especially
important across English-using communities: a familiar surface form can invite
a dominant-variety reading even when the immediate episode supports a
different locally conventional function. Conversely, a static cultural
profile can overcorrect, treating a community-level regularity as a rule about
an individual speaker.

Most context-aware systems make a prediction after context injection. This
leaves two distinct failure modes unresolved. First, the available context may
be missing the one field that distinguishes plausible interpretations or may
contain contradictions. Second, the model may claim confidence while failing
to use context in a behaviorally coherent way. A safe runtime system therefore
needs a decision boundary: it should commit only when the prediction survives
candidate-indexed hypothesis-preserving probes and responds as specified to
hypothesis-contrast probes; otherwise it should ask for one discriminating
fact or withhold the interpretation. These tests characterize model behavior
relative to the model-selected candidate; they do not establish that candidate
as the speaker's true intended meaning.

This paper develops that boundary as \cdcv. It is a calibrated
behavioral-consistency gate, not a semantic-truth verifier, a mechanism for
recovering private ground truth, or a context-authentication system. The gate
checks allowed provenance and deterministic conflicts, then estimates whether
a model's selected candidate behaves consistently under a reviewed probe
schema. Human validation and calibrated error estimates remain necessary: a
consistently wrong model can be invariant, and a probe-responsive model can
still select the wrong sense.

The study asks exactly two confirmatory questions:

\begin{description}[leftmargin=0pt,style=nextline]
\item[RQ1---Selective reliability.]
At matched coverage, does counterfactual context verification reduce wrong
committed interpretations compared with always-answer, confidence-only,
self-consistency and standard context-injection baselines?

\item[RQ2---Robustness and repair.]
Does the proposed system remain invariant under hypothesis-preserving probes,
respond as specified under hypothesis-contrast probes, and recover
gold-\clarify{} cases using one targeted clarification?
\end{description}

Upon completion of the sealed evaluation, the intended safe contribution
sentence is:

\begin{quote}
We design and evaluate a schema-constrained inference-time gate that probes a
model-selected intended-meaning candidate with community-validated
hypothesis-preserving and hypothesis-contrast context probes, then uses
calibrated behavioral-consistency evidence for selective commitment, one
targeted clarification, or abstention.
\end{quote}

Here, ``community validated'' applies to the curated benchmark instances and
reviewed intervention operators. It must not be extended to arbitrary model-
generated deployment instances without a separate fidelity check. The
completed work is intended to contribute (i) an operational dual-probe
behavioral-consistency mechanism, (ii) an action-labelled and expression-family-held-out
three-community English-use benchmark, (iii) a one-question repair policy, and (iv) a paired
equal-budget evaluation. These form a combination claim, not a claim to have
invented selective prediction, context perturbation, clarification, or model
routing individually.

At the present protocol stage, the accurate claim is narrower: we design the
mechanism and specify its evaluation; we do not yet claim improvement in any
outcome.

\section{Background and related work}

\subsection{Pragmatic meaning, variation, and repair}

Pragmatic meaning is an inference about situated communicative action rather
than a context-free dictionary lookup. Contextualization cues connect an
utterance to an activity and to participants' assumptions
\citep{gumperz1982discourse}; common ground is updated across a sequence of
joint actions \citep{clark1996using}. When understanding becomes problematic,
conversation provides organized practices for repair
\citep{schegloff1977repair}. Voice-interface studies likewise show that system
use is sequentially embedded rather than a series of isolated commands
\citep{porcheron2018voice}.

Computational benchmarks now cover multiple pragmatic phenomena
\citep{sravanthi-etal-2024-pub}, while dialect benchmarks document broad
variation in NLP behavior \citep{ziems-etal-2023-multi,
faisal-etal-2024-dialectbench}. Such work supports per-variety evaluation but
does not justify demographic profiling. Dialect cues can also trigger harmful
stereotypes or dominant-variety defaults \citep{fleisig-etal-2024-linguistic,
hofmann2024dialect}. We consequently treat community labels as bounded
resource indexes and self-described language histories, not as causal or
essential properties of speakers. The data documentation and generalization
boundary follow the principles of data statements
\citep{bender-friedman-2018-data}.

The three community resources are bounded by scholarship on postcolonial
English varieties \citep{schneider2007postcolonial}, English use in Korea
\citep{rudiger2019korean}, and situated Nigerian English pragmatic markers
\citep{unuabonah2021abeg,honkanen2022markers}. These sources motivate
expression discovery and contextual dimensions; they do not supply sealed
labels or license a claim that one form has one meaning for an entire country.

\subsection{Ambiguity detection and clarification}

Language models often fail to recognize or represent multiple plausible
readings \citep{liu-etal-2023-afraid}. Prior systems detect ambiguity, ask a
clarification, and answer after the response
\citep{lee-etal-2023-asking}. Repetition-based confidence is a strong baseline
for selective answering of ambiguous questions
\citep{cole-etal-2023-selectively}, and entropy over intent predictions can
identify cases in which interaction helps
\citep{zhang-choi-2025-clarify}. Clarification-question ranking by expected
value of information predates current LLMs
\citep{rao-daume-iii-2018-learning,white-etal-2021-open}.

The closest current work uses structured uncertainty and cost-penalized
expected value of perfect information to decide what to ask and when to stop
in tool-calling agents \citep{suri-etal-2026-structured}. Human and model
clarification behavior also differs by ambiguity type
\citep{madge-etal-2025-referential}. Thus neither structured uncertainty nor
information-value clarification is claimed as new here. Our question is
whether a one-slot clarification policy can repair cases first rejected by a
dialect-context verifier without creating excessive or identity-seeking
questions.

\subsection{Behavioral and counterfactual tests}

Behavioral testing distinguishes invariance expectations from directional
expectations \citep{ribeiro-etal-2020-beyond}. Contrast sets and
counterfactually augmented data use small, meaningful changes that alter a
validated label \citep{gardner-etal-2020-evaluating,
kaushik2020counterfactual}. Recent pragmatic benchmarks already hold an
utterance fixed while varying context, including manually verified context
flips and controlled implicature contexts
\citep{li-etal-2026-born,arai-ren-2026-drinq}. Context flips are therefore not
a standalone novelty.

The present mechanism uses both directions before one runtime decision:
preserving variants test whether incidental context changes destabilize a
prediction, whereas targeted changing variants test whether a minimally
different discourse configuration licenses the reviewed alternate sense. Raw
change is not rewarded; the model must move toward the validated target.

\subsection{Selective prediction, calibration, and routing}

Selective classification trades coverage for lower prediction risk
\citep{geifman2019selectivenet}. Semantic entropy and rejection curves provide
strong LLM uncertainty baselines \citep{farquhar2024semantic}. Conformal risk
control provides a framework for calibrating bounded monotone losses under
stated assumptions \citep{angelopoulos2024conformal}. We use calibrated
selective prediction but claim no finite-sample guarantee unless the selected
procedure's exchangeability and family-level conditions are met.

Quality- and cost-aware routing between models is also established
\citep{ding2024hybrid,ong2025routellm}. In \cdcv, routing is a bounded internal
operation after unresolved verification, not a new output class or a presumed
improvement. A routed prediction must pass the same gate, and degradations are
logged.

\subsection{Research lineage and remaining gap}

Three prior systems provide directly relevant engineering baselines. IUUY described and
prototyped a four-role interaction workflow and a hybrid dialect-table plus
Gemini direction for speech ingestion, interpretation, user-facing
clarification, and feedback provenance \citep{mba2025iuuy}. IEDI introduced
the IEDID resource and formalized a thresholded fuzzy-retrieval path plus a
proposed LLM fallback for three English-use settings \citep{mba2025iedi}. The
frozen public notebook implements only global RapidFuzz matching at threshold
75 and returns \texttt{Unknown} below threshold; it does not implement the
paper's proposed LLM fallback \citep{iedinotebook}. KICS-W later described
contextual persona/codebook prompting, \texttt{Tone\_Category} and
\texttt{Linguistic\_Context} metadata, and a Gemini Pro--Flash dynamic-model-
manager concept \citep{mba2026persona}. These published systems and the
anonymized executable artifact are treated as foundations or reconstruction
baselines, not contributions of the present paper.

The cited literatures motivate testing a narrower combination: validated
expression families drawn from three bounded English-use community resources;
typed, ephemeral context cards;
paired preserving and minimal sense-changing interventions; conflict and
leakage audits; and one calibrated \commit/\clarify/\abstain controller under
family holdout and equal compute. This is a bounded gap statement based on the
cited audit, not an absolute first-ever or patentability claim.

\section{Problem formulation and threat model}

\subsection{Episode and candidate meanings}

Each episode has physically separate runtime and evaluator views:
\begin{align}
r_e&=(x,c,\mathcal{M},\mathcal{T},\mathcal{Q}),\\
g_e&=(y,a^*,J^*).
\end{align}
Here $x$ is the utterance with short co-text; $c$ is an interaction-scoped
context card; $\mathcal{M}=\{m_1,\ldots,m_J\}$ is the frozen candidate set;
$\mathcal{T}$ is a symmetric candidate-indexed probe library; and
$\mathcal{Q}$ is the approved question bank. The sealed evaluator view holds
the community-validated target sense $y$ when recoverable, reference action
$a^*$, and acceptable clarification slots $J^*$. The inference process can
mount $r_e$ but never $g_e$. The action space is
\begin{equation}
\mathcal{A}=\{\commit(m),\ \clarify(j),\ \abstain(r)\},
\end{equation}
where $j$ is a permitted context slot and $r$ is a reason code.

The primary task supplies the same three candidates to every system: two
validated competing senses and an \textsc{Other/Unlisted} option. This isolates
selection and action policy from candidate recall. A secondary end-to-end
condition generates candidates before scoring and reports candidate recall
separately. If the validated sense is absent from $\mathcal{M}$, the event is a
candidate-generation failure; it cannot be repaired by selector calibration.
Selecting \textsc{Other/Unlisted} therefore forces candidate-coverage
abstention and can never produce \commit.

\subsection{Reference actions}

The benchmark separates interpretation from action appropriateness:

\begin{table}[t]
\centering
\caption{Episode evidence state and reference system action.}
\label{tab:actions}
\begin{tabularx}{\linewidth}{>{\raggedright\arraybackslash}p{0.31\linewidth}
>{\raggedright\arraybackslash}p{0.24\linewidth}X}
\toprule
Evidence state & Reference action & Scoring rule \\
\midrule
Sufficient, coherent context & \commit(target) & A sense commitment is correct only if it equals the validated target. \\
Alternate coherent context & \commit(contrast) & The alternate context licenses the independently validated contrasting sense. \\
Discriminating field missing & \clarify(slot) & The slot must belong to the accepted question set. \\
Contradictory, unsafe, or out of scope & \abstain(reason) & No sense commitment is permitted. \\
\bottomrule
\end{tabularx}
\end{table}

\subsection{Threat model}

We target five observable failure classes:

\begin{enumerate}[leftmargin=*]
\item \textbf{Unsupported commitment:} the model selects a sense when context
is missing, contradictory, or out of scope.
\item \textbf{Context neglect:} the prediction remains fixed when a validated
minimal change supports a contrastive sense.
\item \textbf{Context overreaction:} a hypothesis-preserving probe changes the
selected candidate.
\item \textbf{Profile overreach:} a generalized variety cue dominates the
episode or is treated as a deterministic identity rule.
\item \textbf{Evaluation leakage:} test expressions, glosses, contexts, or
labels occur in retrieval resources, prompts, profiles, or study-controlled
development data.
\end{enumerate}

The mechanism does not defend against a deceptive user, compromised model
provider, undetectable pretraining contamination, or arbitrary online
counterfactual generation. Controllable leakage is an offline protocol audit,
not a learned controller feature; otherwise the controller could learn to
recognize benchmark artifacts.

\section{Counterfactual Dialect-Context Verification and Clarification Gate}

\subsection{Runtime flow}

The mechanism is a pre-commit wrapper around a frozen scorer:

\begin{center}
\small
\begin{tabular}{c@{ $\rightarrow$ }c@{ $\rightarrow$ }c@{ $\rightarrow$ }c}
utterance + candidates & typed context card & candidate-indexed probes & calibrated gate\\
&&& $\downarrow$\\[-0.2em]
&&& \commit\quad\clarify\quad\abstain
\end{tabular}
\end{center}

For the curated benchmark, every intervention instance will receive
independent human validation. In a deployment prototype, the runtime may instantiate only a
reviewed operator library. A generated instance that fails schema, fidelity,
or conflict checks forces \abstain; model-generated context is never assumed
valid because the same model can score it.

\subsection{Episode-specific context card}

The context card contains only fields that can alter the local linguistic
inference:
\begin{equation}
c=(v,h,g,r,s,\pi),
\end{equation}
where $v$ is a voluntarily self-declared or experimentally supplied variety
cue, $h$ is the preceding speech act/dialogue evidence, $g$ is the discourse
goal, $r$ is the interaction role or relationship, $s$ is the
situation/register, and $\pi$ records field-level provenance, missingness,
confidence, and expiry. The card is limited to the current interaction.

Race, ethnicity, nationality, religion, gender, sexuality, disability, age,
biometrics, and other protected attributes are neither inferred nor
counterfactually changed. The variety cue remains fixed across intervention
pairs. Relationship \emph{roles} may be manipulated (for example, service
encounter versus peer collaboration) only when the role, rather than personal
identity, is the linguistically relevant variable.

\subsection{Base scoring}

In the end-to-end system, a frozen generator first proposes
\begin{equation}
\mathcal{M}=G_\phi(x)=\{m_1,m_2,\text{\textsc{Other/Unlisted}}\}
\end{equation}
from the utterance and dialogue co-text before context scoring. It cannot read
the sealed label. The confirmatory primary condition replaces $G_\phi$ with
the validated fixed set to isolate context-conditioned selection and gating.
The secondary end-to-end condition restores $G_\phi$ and reports candidate
recall, selection error conditional on recall, and total pipeline error. Thus
the mechanism includes candidate generation without implying that the primary
oracle-candidate experiment solves open-set sense discovery.

The frozen scorer returns
\begin{equation}
p_0(m)=p_\theta(m\mid x,c,\mathcal{M}),\qquad
\hat y=\arg\max_{m\in\mathcal{M}}p_0(m).
\end{equation}
The top-two margin and normalized entropy are
\begin{align}
\Delta_0 &=p_0(\hat y)-\max_{m\ne\hat y}p_0(m),\\
H_0 &=-\frac{\sum_m p_0(m)\log p_0(m)}{\log|\mathcal{M}|}.
\end{align}
If the provider exposes no usable probabilities, a constrained ranking
adapter produces scores and is calibrated separately. Free-form verbal
confidence is not treated as probability without validation.

\subsection{Hypothesis-preserving and hypothesis-contrast probe sets}

The probe library is symmetric across the two non-
\textsc{Other/Unlisted} candidates and is constructed without $g_e$. For each
candidate $m$, define
\begin{align}
\mathcal{C}^{+}_e(m)
 &=\{T_i^+(c,m):d(c,T_i^+(c,m))\leq\delta\},\\
\mathcal{C}^{-}_e(m)
 &=\{(T_j^-(c,m,t_j),t_j):d(c,T_j^-(c,m,t_j))\leq\delta,
 \;t_j\ne m\}.
\end{align}
After base scoring, the gate selects $\mathcal{C}^{+}_e(\hat y)$ and
$\mathcal{C}^{-}_e(\hat y)$ using its own prediction. An operator contract
states that $T_i^+$ is a community-validated \emph{hypothesis-preserving
probe} expected to retain the branch candidate and $T_j^-$ is a
community-validated \emph{hypothesis-contrast probe} expected to support
$t_j$; it does not state which sense is correct in the base episode. Thus a
branch can test candidate-conditioned behavioral consistency but cannot verify
semantic truth, the sealed intended meaning, or the authenticity of context.
The scorer receives only the patched context and common candidate list. It
does not receive relation names, target identifiers, validator scores, or
sealed labels; the controller compares the returned distribution with the
candidate-indexed operator contract afterward.
The utterance, candidate set, and variety cue stay fixed. Changes are confined
to approved slots such as discourse goal, preceding speech act, relationship
role, setting, or situation. Hypothesis-contrast probes normally edit one
discriminating slot and at most two. Candidate glosses, obvious label words,
and identity proxies are prohibited.

A frozen cross-record validator must run before inference. It confirms that
the two distinct branch identifiers match exactly the two runtime
non-\textsc{Other/Unlisted} candidates; every contrast target is the other
core candidate; source context hashes match the base card; no probe references
itself; each patch changes only declared permitted slots; variety and all
fixed fields remain unchanged; and every branch and probe passed the
prespecified acceptance rules. Failure makes the episode ineligible for
commitment and is logged; no replacement branch may be chosen after scores are
observed.

For each $c_i^+$ and $c_j^-$, the same scorer produces distributions $p_i^+$
and $p_j^-$. We compute label invariance and probability retention:
\begin{align}
I^+ &=\frac{1}{K_+}\sum_i
\mathbf{1}[\arg\max_m p_i^+(m)=\hat y],\\
P^+ &=\frac{1}{K_+}\sum_i p_i^+(\hat y),\\
D^+ &=\frac{1}{K_+}\sum_i \operatorname{JSD}(p_0,p_i^+).
\end{align}
JSD uses base-2 logarithms and is bounded by one for the discrete
distributions here \citep{lin1991divergence}.
Targeted responsiveness and contrast margin are
\begin{align}
F^- &=\frac{1}{K_-}\sum_j
\mathbf{1}[\arg\max_m p_j^-(m)=t_j],\\
M^- &=\frac{1}{K_-}\sum_j
\left[p_j^-(t_j)-p_j^-(\hat y)\right].
\end{align}
Validator-agreement weights are a prespecified sensitivity analysis. A model
receives no credit for an arbitrary flip.

\subsection{Verification vector and calibrated controller}

The controller receives the prespecified vector
\begin{equation}
\mathbf{z}_e=[\Delta_0,1-H_0,I^+,P^+,1-D^+,F^-,M^-,V,U,\Gamma],
\end{equation}
where $V$ summarizes intervention validity, $U$ is relevant-slot
completeness, and $\Gamma$ is a deterministic context-conflict score. Budget
is an action constraint, not a correctness feature. Benchmark-leakage flags
remain audit metadata and are not inputs.

Given the small development split, the primary implementation reduces the
features to four monotone groups---base confidence, preservation, targeted
change, and context quality. Its form and weights are chosen/fitted only on
development. A one-dimensional isotonic mapping fitted on the 60-family
calibration split, with equal total weight per family, estimates
\citep{zadrozny2002transforming}
\begin{equation}
q_e\approx P\{a^*=\commit(\hat y)\mid\mathbf{z}_e\},
\end{equation}
the probability that committing is both action appropriate and sense correct.
The score form, monotonicity constraints, calibration procedure, and threshold
$\tau$ are frozen before test access. A conformal-risk controller is a baseline
or sensitivity analysis unless its family-level exchangeability conditions
and loss definition support the claimed guarantee.

Hard policy checks override confidence. A frozen eligibility mask marks an
episode ineligible for commitment after an invalid card or probe bundle,
prohibited field, unresolved hard contradiction, unsafe request,
\textsc{Other/Unlisted} prediction, missing required discriminator,
controllable leakage violation, out-of-domain input, disallowed provenance,
or exhausted hard budget. A high $q_e$ cannot override the mask. Subject to
those constraints,
\begin{equation}
\pi(e)=
\begin{cases}
\commit(\hat y), & q_e\ge\tau,\; V\ge v_{\min},\;\Gamma=0,\\
\clarify(j^*), & \text{one permitted discriminating slot is unresolved and }J(q^*)\ge\kappa,\\
\abstain(r), & \text{otherwise.}
\end{cases}
\end{equation}

\subsection{One targeted clarification}

The question bank contains neutral, answerable, community-reviewed templates
mapped to one permitted slot. Let $\mathcal{Q}_e$ be the valid questions for
the episode. The gate selects
\begin{align}
q^*=\arg\max_{q\in\mathcal{Q}_e}\;&
\left[H(p_0)-\sum_{a\in\mathcal{A}_q}
\widehat P(a\mid e,q)H\{p_\theta(\cdot\mid x,c\oplus(q,a))\}\right]
\nonumber\\
&-\lambda_I C_I(q)-\lambda_P C_P(q)-\lambda_C C_C(q),
\end{align}
where the terms penalize interaction burden, privacy sensitivity, and compute.
$C_P(q)=\infty$ for a prohibited field. Answer priors and penalty weights come
only from development/calibration data. The question must not reveal a
candidate gloss or presuppose a community stereotype.

The primary selector reads a frozen clarification-scenario manifest. For every
approved case--question--answer scenario, the manifest records a normalized
answer prior, a reusable probe-score or hypothetical-score reference, its
source split and immutable hash, and the validated context-patch reference.
The runtime view contains no reference action or sense. In the primary
implementation, every score reference resolves to a distribution already
produced within the allocated probe pass or to a development/calibration-only
static table. Question ranking therefore adds zero sealed-test model calls.
A missing or invalid scenario entry makes the question ineligible. Exploratory
live hypothetical scoring is reported separately and charged to both the
system and comparator budgets.

For the controlled RQ2 experiment, an access-controlled harness releases one
independently validated answer patch only after the system emits the matching
question for that case and slot. A wrong or unnecessary question receives no
helpful patch. The answer is not generated by a model and is not a live-user
observation. The card records provenance as
\texttt{standardized\_clarification}; the full verifier runs once more and
must commit, route once, or abstain/escalate. A second question is prohibited.
Clarification recovery therefore measures controlled repair potential, not
natural user behavior.

\subsection{Bounded larger-model routing}

If context is complete and coherent but model uncertainty remains, a
development-fitted router may estimate the larger model's incremental benefit.
It routes only when expected benefit exceeds monetary/latency cost and privacy
policy permits transmission. The routed answer passes the same verifier. We
report helpful and harmful routes separately; a larger model is not presumed
to be safer. Routing is inherited engineering, not the paper's novelty.

With $K=K_++K_-$ interventions and scorer cost $T_f$, one verification pass is
$O((1+K)T_f+K|\mathcal{M}|)$ after candidate construction. Features can be
streamed in $O(|\mathcal{M}|)$ memory. In the confirmatory budget,
$K_+=K_-=1$, so one pass makes three scoring calls. The primary question
selector adds zero calls. All comparisons receive both equal maximum calls and
equal mean token budgets. With at most one initial pass, one post-question
pass, and one routed pass, the cap is nine scorer calls excluding candidate
generation. A live hypothetical-call variant is exploratory and receives a
larger, explicitly matched budget rather than being hidden inside this cap.

\section{Counterfactual three-community English-use benchmark}

\subsection{Role of current IEDID rows and personas}

Current IEDID rows and inherited personas can be used for expression
discovery, development examples, context-card schema construction, and
lineage baselines. They cannot constitute sealed-test gold. The current
resources were not designed with paired interventions, action labels, creator-
validator separation, or reliable phrase/gloss leakage boundaries. Some live
records also require deidentification, provenance, consent, and license review.

Before use, the development snapshot is pinned, sanitized, and assigned row-
level provenance. No test expression, sense gloss, intervention, question, or
action label is sampled directly from it. The previously audited revision is
recorded for provenance, but the final experimental revision remains unfrozen
until governance checks pass.

\subsection{Family design and split}

An expression family is the independent unit. Each retained family has two
validated contrasting senses and eight nested episode cases:

\begin{enumerate}[leftmargin=*]
\item canonical context for sense A;
\item canonical context for sense B;
\item one meaning-preserving variant of A;
\item one meaning-preserving variant of B;
\item one minimal meaning-changing variant from A toward B;
\item one minimal meaning-changing variant from B toward A;
\item one underspecified case labelled \clarify; and
\item one contradictory or out-of-scope case labelled \abstain.
\end{enumerate}

For every runtime episode, a probe team separately constructs one
hypothesis-preserving and one hypothesis-contrast operator for each of the two
core candidate branches without seeing the sealed base action or sense. Thus
both branches exist even when one is inconsistent with the original context.
The model chooses a branch only from $\hat y$. Probe validators assess the
candidate-relative operator contract but not the base reference label. These
are community-validated candidate-indexed probes, not claims about the
episode's unknown true meaning. No operator may insert a candidate gloss or
obvious label word.

The recommended retained sample is 30 development, 60 calibration, and 150
sealed-test families. The test contains 50 families for each of the three
community resources and 1,200 cases. The inferential sample is 150 families,
not 1,200 independent observations. A 5,000-replication planning simulation
supported 150 families for a pooled nominal five-percentage-point paired
effect under its base assumptions. That simulation does not establish an
effect, does not power within-community superiority, and supplies no result to
this paper. At 70\% expected validation retention and balanced recruitment,
approximately 345 candidate families (115 per community) are needed to retain
240; a blinded validation pilot will update the recruitment target.

The retained design contains exactly 1,920 base episode cases
($240\times8$). Four symmetric probe instances are authored per base case
(two candidate branches times one hypothesis-preserving and one
hypothesis-contrast probe), yielding 7,680 probe instances and 38,400 probe
ratings at five validators each. Base-case validation contributes 9,600
ratings. Because there is one gold-\clarify{} case per family, the retained
splits contain 240 such cases (150 in the sealed test). One accepted question
per case would require 240 items and 1,200 ratings; the planned
information-gain-versus-random comparison requires at least two accepted
options per case, hence at least 480 accepted items and 2,400 ratings. If
$Q$ candidate question items, including rejected items, are submitted, the
actual question workload is $5Q$ with $Q\ge480$. The minimum accepted-item
workload is therefore 50,400 ratings (9,600 base + 38,400 probe + 2,400
question ratings), excluding rejected candidates, qualification tasks, and
adjudication. The final $Q$, replacement allowance, completion-time estimate,
compensation, platform fees, and adjudication budget are frozen before ethics
approval and recruitment.

\subsection{Community authorship and validation}

Creators and validators are recruited through self-described relevant
language histories and community familiarity, not nationality alone. A
creator cannot validate the same bundle. Five independent variety-proficient
validators first provide an open interpretation and then assess the candidate
choice, system action, naturalness, coherence, minimality, leadingness, and
sensitivity. Acceptance requires at least four of five target/action choices;
median naturalness and coherence of at least four on a five-point scale;
median minimality of at least four for meaning-changing edits; no detected
gloss leakage; and no unresolved safety or provenance flag. Accepted questions
also require median leadingness at most two (one means not leading).
The governance process adapts participatory research principles while keeping
decision authority and limitations explicit
\citep{nekoto-etal-2020-participatory}.

Agreement is reported as raw distributions, entropy, and Krippendorff's alpha
where the scale permits \citep{artstein2008intercoder}. Consensus is an
operational target label, not direct access to a speaker's private intention
or evidence that every community member shares the judgment. Community
partners adjudicate uncertain cases and review quoted examples and failure
interpretations before release.
Consent distinguishes expression authorship, base-case rating,
candidate-indexed probe rating, clarification-question rating, possible
quotation, retention, controlled release, and withdrawal. Recruitment proceeds
only after the ethics and funding reviews confirm that the full workload,
community balance, compensation, privacy controls, and rejected-item
replacement allowance are feasible. If the design must be resized, power is
resimulated and the protocol is updated before preregistration.

\subsection{Leakage separation and sealing}

Splits are expression-family disjoint. Before sealing, normalized expressions,
paraphrases, sense glosses, profile phrases, prompts, and retrieval indexes are
hashed and compared. The audit covers every IEDID snapshot used by a system,
all inherited paper examples and profiles, few-shot prompts, development and
calibration families, and clarification templates. Lexical and embedding
alerts receive blinded human adjudication. Any controllable overlap is removed
before test hashes are published.

Proprietary foundation-model pretraining cannot be fully audited. The study
therefore reports this residual limitation, runs memorization-oriented probes,
and avoids claiming proof of contamination absence. Leakage flags are used to
invalidate or qualify cases, never to improve the learned gate.

The inference job mounts only a schema-valid runtime view, probe bundles,
context cards, and question templates. It cannot mount reference actions,
senses, case types, acceptable slots, or post-answer reference senses.
Predictions and resource logs are hashed and frozen before a separate evaluator
joins the access-controlled labels by case identifier. This separation is
enforced in schema and package tests, not only stated in prose.

\section{Experimental design}

\subsection{Systems and baselines}

All systems receive identical candidate ordering seeds and frozen inputs
allowed by their condition. The comparison set is:

\begin{enumerate}[leftmargin=*]
\item \textbf{IEDI-NB:} frozen public-notebook fuzzy lookup, reported as a
lineage/coverage diagnostic rather than a competitive engineering baseline;
\item \textbf{Frozen contextual ranker:} a non-LLM bi-encoder or cross-encoder
that scores the utterance, typed card, and candidate definition;
\item \textbf{Direct LLM:} utterance and dialogue co-text, always answer;
\item \textbf{Direct-Random@$C^*$:} a label-blind case-hash subsample of the
always-answer commitments, included only to project that baseline to common
coverage without using confidence;
\item \textbf{KICS-W\_RECON:} paper-faithful persona/codebook reconstruction,
with all prompt/codebook and router assumptions disclosed;
\item \textbf{Structured context:} typed card without intervention probes,
always answer;
\item \textbf{Score gate:} token/ranking margin and normalized entropy as
separate frozen scores;
\item \textbf{Sample-agreement gate:} Cole-style repeated sampling and modal
agreement;
\item \textbf{Semantic-entropy gate:} meaning-clustered sampled outputs;
\item \textbf{Calibrated gate:} isotonic/logistic mapping of the structured-
context score;
\item \textbf{Conformal-risk control:} a separate implementation with stated
family-level unit, bounded loss, and exchangeability assumptions;
\item \textbf{Probe-only aggregation:} the same symmetric probe library and
budget as \cdcv, but a simple frozen aggregation without the joint controller;
\item \textbf{Structured-uncertainty question gate:} the same question bank,
answer harness, and cost penalties but no dual counterfactual verification;
\item \textbf{\cdcv{}--no question:} full verification, commitment or
abstention only;
\item \textbf{\cdcv{}+one question:} full proposed repair policy; and
\item \textbf{Equal-budget context control:} structured context with the same
maximum calls and mean token allowance as \cdcv.
\end{enumerate}

The equal-budget context control is the primary comparator. It makes three
independent stochastic calls with the identical utterance and structured card,
uses the mean candidate distribution plus sample agreement, and fits the same
one-dimensional isotonic calibration on the same family-weighted calibration
split; it receives no intervention relation. Temperature, seeds, aggregation,
maximum calls, and mean token cap are frozen. Dialogue-only sample agreement
and semantic entropy receive the same call/token budgets because extra
inference alone can improve reliability. Allocated and consumed calls/tokens
are both reported.
Model identifiers, provider dates, decoding, prompts, candidate order,
hardware, batch size, cache policy, and retry rules are frozen.

Always-clarify and oracle-slot-question policies are reported only as burden
and controlled-repair bounds. A random validated question is the negative
control. They are not competitive systems.
The native always-answer result remains at coverage one; Direct-Random@$C^*$
is its explicitly non-deployable matched-coverage projection. The confidence-
only gate is the deployable selective version.

\subsection{IEDI protocol correction}

If every sealed family is removed from the frozen IEDI lexicon, IEDI-NB is
expected to return \texttt{Unknown}. We report that regime only as a lineage
coverage diagnostic. A separately labelled secondary regime may test
\textbf{known expression, unseen context}: the expression-to-entry mapping is
frozen, but every episode context, intervention, question, and action label is
new. The two regimes are never pooled.

\subsection{Prespecified bounded application pilot}

A separately labelled, non-confirmatory pilot may begin only after the sealed
benchmark analysis is complete and a pilot-specific sample-size or precision
plan, interface build, models, prompts, thresholds, outcomes, and stopping
rules have been frozen. The pilot embeds \cdcv{}+one question and the
structured-uncertainty question comparator in a non-high-stakes, text-only
mediation sandbox. It uses newly consented interactions that are disjoint from
development, calibration, and confirmatory test families; pilot records cannot
tune or replace any confirmatory component.

Participants supply or confirm the intended interpretation before system
feedback and confirm or reject the final interpretation after at most one
clarification. Prespecified pilot outcomes are participant-confirmed
wrong-commit burden, appropriate commitment or abstention, one-question repair,
unnecessary questions, abandonment, interaction burden, latency, tokens,
calls, memory, and cost. Recruitment requires separate ethics review and
consent for participation, retention, quotation, and release. The interface
collects no audio, infers no protected attribute, creates no persistent
cultural profile, and triggers no high-stakes action. Analysis is limited to
the frozen pilot setting with participant- and expression-family-clustered
uncertainty; it supports no confirmatory superiority or real-world outcome
claim.

\subsection{Metrics}

Let $N$ be the number of scored episodes, $C$ the set of committed episodes,
and $W\subseteq C$ the wrong commitments. A commitment belongs to $W$ when
the reference action is noncommit or when a commit-eligible episode receives
the wrong sense. We report
\begin{align}
\text{coverage} &= |C|/N,\\
\text{selective risk} &= |W|/|C|,\\
\text{overall wrong-commit burden} &= |W|/N.
\end{align}
The primary outcome is selective risk at matched coverage over all eight case
types; safe action-appropriate coverage is capped at $6/8=0.75$. The overall burden
prevents conditional reporting from hiding abstention volume. We also report
interpretation-only risk on the six commit-eligible cases and unsupported-
commit rates on clarify/abstain cases. The full risk--coverage curve is
descriptive. Over the prespecified common range $[C_{\min},C_{\max}]$,
\begin{equation}
\operatorname{AURC}=\frac{1}{C_{\max}-C_{\min}}
\int_{C_{\min}}^{C_{\max}}R(C)\,dC,
\end{equation}
with lower values better; the trapezoidal, zero-commit, and tie rules are
frozen in code.

Required action and robustness metrics are the three-way action confusion
matrix, action macro-F1, preservation invariance, preservation correctness,
appropriate target-flip rate, flip coverage, clarification recovery,
unnecessary-question rate, and final action-aware accuracy after one question.
Candidate recall is reported in the generated-candidate condition. Calibration
metrics include Brier score and expected calibration error
\citep{brier1950verification,guo2017calibration}.

Hypothesis-preservation denominators include every accepted
hypothesis-preserving probe for an eligible non-Other base branch;
preservation correctness additionally requires
the sealed sense to remain correct. Targeted-response denominators include
every accepted hypothesis-contrast probe actually run, and flip coverage reports the
fraction of episodes with a valid non-Other branch. Clarification recovery is
computed primarily over all gold-\clarify{} cases in the evaluated split
($N=150$ in the sealed test). A success requires an acceptable matching
question followed by a correct commit after the standardized answer; not
asking, asking a nonmatching question, or failing to commit correctly counts
as failure. Recovery conditional on actually asking a matching question is a
secondary diagnostic and is accompanied by the matching-question rate.
Unnecessary-question
rate is the proportion of sufficient or abstain-labelled cases on which a
question is asked. Final action-aware accuracy requires both the correct final
action and, for commitments, the correct sense.

Every metric is disaggregated by community resource with confidence intervals
and a prespecified harm guard. Runtime reporting includes p50/p95 latency,
input/output tokens, calls, peak memory, and cost under a dated price snapshot.
Routing frequency, beneficial routes, and harmful routes are included. Failure
analysis codes dominant-variety default, profile overreach, context neglect,
context overreaction, conflict miss, unsafe question, invalid output, candidate
omission, and controllable leakage.

\subsection{Primary estimand and inference}

RQ1 is evaluated before user clarification to isolate verification. The
primary contrast is \cdcv{}--no question versus the equal-budget structured-
context control. A common target coverage $C^*\leq0.75$ is selected on the
calibration split, is no greater than the minimum eligible calibration
coverage of the two primary methods, and is then frozen. Each method applies
its hard eligibility mask before ranking. On test, it ranks only eligible
cases under the frozen score rule without reading labels; a frozen case hash
resolves boundary ties. Invalid cards or bundles, prohibited fields, hard
conflicts, \textsc{Other/Unlisted}, missing discriminators, controllable
leakage violations, out-of-scope inputs, and hard-budget failures are never
forced to \commit. If fewer than $C^*N$ test cases are eligible, the matched
target is reported as unattainable for that method, together with its observed
coverage and risk. At attainable matched coverage we estimate
\begin{equation}
\Delta_R=R_{\mathrm{CDCV}}-R_{\mathrm{control}}.
\end{equation}
We use 10,000 paired expression-family bootstrap resamples, stratified by
community, with method-specific eligibility masks preserved and score ranking
recomputed only among eligible cases inside each resample. A replicate with an
eligibility shortfall is flagged rather than repaired by forced commitment.
The primary criterion is that the upper endpoint of the two-sided
95\% confidence interval for $\Delta_R$ is below zero. At the separately
calibration-frozen deployed threshold, proposed coverage may not be more than
five percentage points below the control.

RQ2 reports hypothesis preservation, targeted contrast response, and
one-question repair. Its primary repair denominator is all gold-\clarify{}
cases, not each policy's selectively rejected subset; conditional recovery is
secondary. Holm
correction is applied to four prespecified secondary contrasts: full versus
no-preserving probes on preservation correctness; full versus no-changing
probes on targeted response; \cdcv{}+question versus the structured-
uncertainty question gate on clarification recovery; and information-gain
versus random validated question on final action-aware accuracy.
The sequential adjustment controls the familywise error rate
\citep{holm1979simple}.
Hierarchical/mixed-effects models are sensitivity analyses with
cases nested in families; they do not replace the paired primary estimand.
The 50-family community strata support heterogeneity and harm-guard reporting,
not independent claims of superiority. A community is flagged when the upper
95\% confidence endpoint permits a wrong-commit-risk increase greater than
five percentage points relative to the control.

\subsection{Ablations and negative controls}

Mechanism ablations remove preserving probes, changing probes, conflict gates,
calibration, clarification, and routing one at a time. A random validated
question tests the information-value selector. Always-large routing measures
whether gate-triggered routing adds value. Fixed versus generated candidates
separate selection from recall. Mismatched context cards/personas,
mismatched-family probes, and target-label-swapped interventions are negative
controls for profile overreach, gloss leakage, and circularity.

The same model must not generate and approve its own test interventions.
Benchmark instances are independently authored and human validated. Any
secondary online generator is frozen, uses reviewed operators, and is scored
for fidelity separately.

\subsection{Calibration and run lock}

Development families select prompt wording, controller form, question costs,
and operational caps. Calibration families fit score mappings and thresholds.
Before test execution we freeze code and environment, split/data hashes,
models and dates, prompts, random seeds, token/call budgets, estimands,
confidence-interval code, and empty table shells. Each system is then executed
once against sealed labels. A post-unsealing correction requires a dated
deviation record and new run identifier.

\section{Results}

\resultlockbox

Table~\ref{tab:rq1} is the confirmatory RQ1 shell. Every cell must be generated
from the immutable sealed run; development, calibration, and simulated values
are prohibited.

\begin{table}[t]
\centering
\caption{RQ1 selective reliability on 150 sealed expression families. Values
will include family-clustered 95\% confidence intervals.}
\label{tab:rq1}
\small
\begin{tabular}{lcccc}
\toprule
System & Coverage & Selective risk & Overall wrong commit & AURC \\
\midrule
Direct LLM & \resultcell & \resultcell & \resultcell & \resultcell \\
Direct-Random@$C^*$ & \resultcell & \resultcell & \resultcell & \resultcell \\
KICS-W\_RECON & \resultcell & \resultcell & \resultcell & \resultcell \\
Frozen contextual ranker & \resultcell & \resultcell & \resultcell & \resultcell \\
Structured context & \resultcell & \resultcell & \resultcell & \resultcell \\
Margin/normalized entropy & \resultcell & \resultcell & \resultcell & \resultcell \\
Sample agreement, equal budget & \resultcell & \resultcell & \resultcell & \resultcell \\
Semantic entropy, equal budget & \resultcell & \resultcell & \resultcell & \resultcell \\
Isotonic/logistic calibrated gate & \resultcell & \resultcell & \resultcell & \resultcell \\
Conformal-risk control & \resultcell & \resultcell & \resultcell & \resultcell \\
Equal-budget context control & \resultcell & \resultcell & \resultcell & \resultcell \\
Probe-only aggregation & \resultcell & \resultcell & \resultcell & \resultcell \\
\cdcv{}--no question & \resultcell & \resultcell & \resultcell & \resultcell \\
\bottomrule
\end{tabular}
\end{table}

\begin{table}[t]
\centering
\caption{Frozen IEDI-NB lineage diagnostic. The regimes remain separate and
are not pooled with the competitive RQ1 table. `Unknown' is mapped to
abstention.}
\label{tab:iedi}
\small
\begin{tabular}{lcccc}
\toprule
Regime & Coverage & Unknown/abstain & Selective risk & Overall wrong commit \\
\midrule
Family held out & \resultcell & \resultcell & \resultcell & \resultcell \\
Known expression, unseen context & \resultcell & \resultcell & \resultcell & \resultcell \\
\bottomrule
\end{tabular}
\end{table}

\begin{table}[t]
\centering
\caption{RQ2 counterfactual robustness before clarification.}
\label{tab:rq2}
\small
\begin{tabularx}{\linewidth}{lXXXX}
\toprule
System & Preservation invariant & Preservation correct & Targeted flip & Flip coverage \\
\midrule
Structured context & \resultcell & \resultcell & \resultcell & \resultcell \\
Sample agreement, equal budget & \resultcell & \resultcell & \resultcell & \resultcell \\
Probe-only aggregation & \resultcell & \resultcell & \resultcell & \resultcell \\
\cdcv{}--no question & \resultcell & \resultcell & \resultcell & \resultcell \\
\bottomrule
\end{tabularx}
\end{table}

\begin{table}[t]
\centering
\caption{RQ2 controlled one-question repair. Oracle-slot is an analysis upper
bound, not a deployable comparator.}
\label{tab:repair}
\small
\begin{tabularx}{\linewidth}{lXXXX}
\toprule
Question policy & Clarification recovery & Unnecessary question & Final action-aware accuracy & Mean questions \\
\midrule
Always clarify & \resultcell & \resultcell & \resultcell & \resultcell \\
Random validated question & \resultcell & \resultcell & \resultcell & \resultcell \\
Oracle slot (upper bound) & \resultcell & \resultcell & \resultcell & \resultcell \\
Structured uncertainty/EVPI & \resultcell & \resultcell & \resultcell & \resultcell \\
\cdcv{}+one question & \resultcell & \resultcell & \resultcell & \resultcell \\
\bottomrule
\end{tabularx}
\end{table}

\begin{table}[t]
\centering
\caption{Measured engineering costs under matched budgets.}
\label{tab:cost}
\small
\begin{tabular}{lcccccc}
\toprule
System & Calls & Tokens & p50 latency & p95 latency & Peak memory & Cost/episode \\
\midrule
Sample agreement & \resultcell & \resultcell & \resultcell & \resultcell & \resultcell & \resultcell \\
Equal-budget context control & \resultcell & \resultcell & \resultcell & \resultcell & \resultcell & \resultcell \\
\cdcv{}--no question & \resultcell & \resultcell & \resultcell & \resultcell & \resultcell & \resultcell \\
\cdcv{}+one question & \resultcell & \resultcell & \resultcell & \resultcell & \resultcell & \resultcell \\
\bottomrule
\end{tabular}
\end{table}

The final Results section will additionally contain the action confusion
matrix, risk--coverage curves, calibration plot, per-community estimates,
ablations, negative controls, known-expression IEDI regime, routing benefit,
and qualitative errors. No directional sentence will be written until the
locked analysis supports it.

\section{Discussion framework}

The completed discussion will interpret a three-way trade-off rather than
headline accuracy alone. A useful gate must lower harmful commitments without
withdrawing service disproportionately from one community. It must also
distinguish two sources of rejection: missing specification, which may be
repairable through one question, and residual model uncertainty or invalid
context, which should lead to abstention or controlled escalation.

Several outcome patterns are informative even if the primary hypothesis
fails. High hypothesis preservation with low correctness would reveal stable error.
Frequent changes with low targeted-flip accuracy would reveal sensitivity
without appropriate context use. Strong repair accompanied by a high
unnecessary-question rate would indicate an intrusive policy. Lower risk
caused solely by reduced coverage would not establish the intended technical
effect. A larger routed model that degrades verified cases would challenge the
assumption that scale is a safe fallback.

The intervention analysis is behavioral, not causal evidence about people.
Changing a discourse goal in a validated vignette tests a model decision
boundary. It does not show that changing a human relationship or community
identity causes a particular meaning in unconstrained interaction.

\section{Limitations}

First, the gate cannot prove that supplied context is factually true,
authenticate its source, or recover a speaker's private ground truth. It can
check declared provenance categories and conflicts and test candidate-relative
model behavior under controlled probes. Second, hypothesis preservation and
targeted contrast responsiveness do not guarantee semantic correctness;
labels and calibration are indispensable.
Third, two primary senses simplify pragmatic continua and multifunctionality.
The \textsc{Other/Unlisted} option and end-to-end candidate-recall analysis
expose but do not solve candidate omission.

Fourth, human validation is bounded by recruited participants, tasks, and
vignettes. Consensus does not represent all speakers. Fifth, counterfactuals
may accidentally lead the model through lexical cues; minimality ratings,
gloss-overlap audits, and label-swapped controls reduce but do not eliminate
this risk. Sixth, proprietary-model training contamination cannot be excluded.
Seventh, calibration may fail under model updates, temporal drift, or a new
community; versions must be frozen and recalibration is required after change.

Eighth, additional calls may make the gate too slow or costly. Equal-budget
controls and measured engineering metrics are therefore part of the central
test. Ninth, a larger model can be worse and routing may create privacy or
cost risk. Finally, three English-using communities do not represent all World
Englishes, and text-based intended-meaning selection does not establish ASR,
prosody, affect, or real-world communication outcomes.

\section{Ethics, privacy, and community governance}

No data collection begins before institutional ethics review or documented
exemption. Consent materials cover storage, quotation, controlled release, and
withdrawal. Creators and validators are compensated, and community partners
review labels, failure interpretations, and example quotations before public
release. Before unsealing, only split and manifest hashes are public. After the
analysis is frozen, releases prioritize the consent-authorized deidentified
structured records required for independent table regeneration; raw audio and
identity-bearing text are excluded by default.

Context is self-declared or directly supported by the current interaction,
expires with the episode, and remains visible and correctable. The system does
not infer protected identities from language. Protected attributes are never
intervention variables or clarification targets. Corrections enter a reviewed,
provenance-preserving queue; they do not automatically rewrite a shared
lexicon or model. High-stakes autonomous decisions and persistent cultural
profiles are outside the deployment scope.

\section{Reproducibility and artifact plan}

The artifact package contains the typed schemas, study configuration, lineage
register, validation script, and empty Results shells. Before evaluation it
will add the ethics-approved annotation instrument, validator qualification
rubric, deidentified aggregate agreement report, split hashes, prompt/model
manifests, controller code, equal-budget scheduler, statistical analysis, and
environment lock. Each prediction log records case, family, model, prompt,
candidate order, controller and intervention hashes together with component
scores, action reason, calls, tokens, latency, memory, and price snapshot.

Subject to ethics, license, consent, provenance, deidentification, and community
review, the post-analysis reproducibility bundle will contain deidentified
cases and reference labels, probe contracts, approved questions and
standardized clarification answers, frozen predictions, model/prompt/controller
manifests, and evaluation code sufficient to regenerate every reported table.
Public release is the default for authorized structured records; controlled
access is used only when the applicable consent, license, or community review
requires it, with every withheld field and reason documented. Identity-bearing
source text and raw audio are never public by default. Hashes and aggregates
support integrity auditing but are not presented as a substitute for the
records needed to reproduce the reported analyses.

\section{Conclusion}

Context-aware interpretation needs a permission-to-commit mechanism. \cdcv
operationalizes that mechanism by testing both invariance under reviewed
candidate-relative hypothesis-preserving probes and specified response to
reviewed hypothesis-contrast probes, then mapping calibrated behavioral-
consistency evidence to commitment, one-slot
clarification, or abstention. The paper's empirical claim remains deliberately
open: only a sealed, family-disjoint, community-validated and equal-budget
evaluation can show whether the gate actually reduces wrong commitments. The
present draft fixes the question, method, comparison, and evidence boundary so
that the eventual answer---positive, null, or adverse---can be reported without
post hoc redesign.

\section*{Declaration of generative AI and AI-assisted technologies}

\textcolor{lockred}{\textbf{AUTHOR ACTION REQUIRED:}} Complete this statement
according to the target journal's policy, identifying any language, code, or
analysis assistance and confirming that the authors verified all content.

\section*{Funding}

\textcolor{lockred}{\textbf{AUTHOR ACTION REQUIRED.}}

\section*{CRediT authorship contribution statement}

\textcolor{lockred}{\textbf{AUTHOR ACTION REQUIRED; do not infer roles from
author order.}}

\section*{Declaration of competing interest}

\textcolor{lockred}{\textbf{AUTHOR ACTION REQUIRED, including any relevant
company affiliation or intellectual-property activity.}}

\section*{Data availability}

The study data do not yet exist. After the sealed analysis, the
consent-authorized deidentified cases, labels, probe contracts, standardized
clarification answers, frozen predictions, and table-regeneration code will be
deposited under a persistent identifier. Records are placed under controlled
access only where consent, license, or community governance requires it. Until
those gates are met, this document remains a protocol/preprint and makes no
completed-data or reproducibility claim.

\bibliographystyle{elsarticle-harv}
\bibliography{references}

\end{document}
